% ============================================================
% MÓDULO 8: PROYECTO FINAL Y FRONTERAS DEL RL
% ============================================================
\section{M8 · Proyecto y Fronteras}
\module{Módulo 8: Proyecto Integrador y Fronteras}{Diagnóstico experimental, diseño de recompensas y el camino hacia el RL moderno}

\begin{frame}{El Proyecto Final: OptiBot en el Almacén Dinámico}
  \begin{columns}[c,onlytextwidth]
    \column{0.50\textwidth}
      \textbf{El Desafío de Integración:}
      \begin{itemize}
        \item Ponemos a prueba a OptiBot en un entorno que cambia en tiempo real:
        \begin{itemize}
          \item Estantes móviles y pasillos bloqueados.
          \item Zonas de batería baja y estaciones de recarga.
          \item Pedidos urgentes con ventanas de entrega estrictas.
        \end{itemize}
        \item \textbf{Evaluación Comparativa:}
        \begin{enumerate}
          \item Política Aleatoria (Línea base)
          \item Q-Learning Tabular (Módulo 6)
          \item Deep Q-Network (Módulo 7)
        \end{enumerate}
      \end{itemize}
    \column{0.48\textwidth}
      \centering
      \begin{tikzpicture}[x=0.075cm,y=0.050cm]
        \draw[-{Latex},SoftText] (0,0) -- (55,0) node[right,font=\tiny] {Episodios};
        \draw[-{Latex},SoftText] (0,-10) -- (0,55) node[above,font=\tiny] {Recompensa media};
        % Aleatorio
        \draw[SoftText,dashed,line width=1pt] (0,-5) -- (50,-5) node[right,font=\tiny] {Aleatorio};
        % Q-Learning Tabular
        \draw[Gold!90!Ink,line width=1.5pt] (0,-8) .. controls (10,5) and (25,25) .. (50,30) node[right,font=\tiny] {Tabular};
        % DQN
        \draw[Teal,line width=2pt] (0,-8) .. controls (15,-5) and (25,35) .. (50,48) node[right,font=\tiny] {DQN};
      \end{tikzpicture}
  \end{columns}
  \vspace{0.3cm}
  \keyidea{El agente DQN demuestra capacidad de adaptación ante configuraciones no vistas durante el entrenamiento.}
\end{frame}

\begin{frame}{Diagnóstico y Patologías Comunes en RL}
  \begin{columns}[T,onlytextwidth]
    \column{0.48\textwidth}
      \begin{alertblock}{¿Por qué no aprende mi agente?}
        \begin{itemize}
          \item \textbf{Exploración insuficiente:} $\varepsilon$ decayó demasiado rápido; el agente nunca descubrió la meta.
          \item \textbf{Recompensas demasiado dispersas (sparse):} Si sólo hay $+100$ al final de 1,000 pasos, el gradiente no fluye.
          \item \textbf{Tasa de aprendizaje inestable:} Pérdida que oscila salvajemente o diverge a $\infty$.
        \end{itemize}
      \end{alertblock}
    \column{0.48\textwidth}
      \begin{block}{El Peligro del Reward Hacking}
        Si recompensamos a OptiBot por ``no colisionar'', el robot puede decidir \textbf{quedarse completamente quieto}.\\
        \vspace{0.2cm}
        \textbf{Lección de diseño:}
        El agente no maximiza lo que tú \emph{deseas} que haga; maximiza estricta y despiadadamente la función matemática que \emph{programaste}.
      \end{block}
  \end{columns}
\end{frame}

\begin{frame}{Fronteras: Del Control Discreto al RL Moderno}
  \begin{columns}[T,onlytextwidth]
    \column{0.32\textwidth}
      \begin{block}{Policy Gradients}
        \textbf{REINFORCE y PPO}\\
        \vspace{0.15cm}
        Optimizan directamente la política $\pi_\theta(a \mid s)$ sin tabla Q. Esencial para controlar brazos robóticos en espacios continuos.
      \end{block}
    \column{0.32\textwidth}
      \begin{alertblock}{Actor-Critic}
        \textbf{A2C, SAC, TD3}\\
        \vspace{0.15cm}
        Combinan un Actor (política que decide) con un Crítico (función de valor que evalúa el error TD). Máxima eficiencia muestral.
      \end{alertblock}
    \column{0.32\textwidth}
      \begin{exampleblock}{RLHF en LLMs}
        \textbf{Alineación de Modelos}\\
        \vspace{0.15cm}
        El aprendizaje por refuerzo con feedback humano es el motor que alinea ChatGPT, Claude y Gemini hacia respuestas útiles y seguras.
      \end{exampleblock}
  \end{columns}
  \vspace{0.4cm}
  \keyidea{Todos estos avances avanzados se fundamentan en los mismos tres pilares que dominaste: \textbf{Probabilidad, Bellman y Bootstrapping}.}
\end{frame}
